\قسمت{گسترش \متن‌لاتین{CSS}}

پرهام به رنگ نارنجی علاقه دارد، او برای طراحی یک وب‌سایت از کتابخانه‌ی
\متن‌لاتین{Bootstrap}
استفاده کرده است. در این کتابخانه برای درست کردن
یک لیست، از کلاس \متن‌لاتین{.list-group-item}
استفاده می‌شود.

به پرهام کمک کنید که بتواند به این کلاس‌ها صفت\پانویس{attribute} رنگ نارنجی را اضافه کند.
در نظر داشته باشید که ما در رابطه با محتوای این کلاس اطلاعی نداریم
پس راه‌حلی ارائه دهید که اگر این کلاس صفات رنگ را هم داشته باشد، باز راه‌حل شما قابل استفاده باشد.
دقت کنید نیاز است کد \متن‌لاتین{CSS} را در کنار توضیحات لازم بیاورید.

ترتیب \متن‌لاتین{import} کردن فایل‌های \متن‌لاتین{css} را مشخص کنید
و استفاده از کلیدواژه \متن‌لاتین{IMPORTANT} مجاز نیست.

\نمره{۲}

% پاسخ پیشنهادی است و پیش از استفاده در تصحیح نیاز به بازبینی دارد.
\begin{پاسخ}

کلاس را در فایلی جدا بازنویسی می‌کنیم و آن فایل را \متن‌سیاه{پس از} \متن‌لاتین{Bootstrap} وارد می‌کنیم:

\begin{latin}
\begin{minted}[]{css}
/* after bootstrap.css */
.list-group-item {
  color: orange;
}
\end{minted}
\end{latin}

اگر خود \متن‌لاتین{Bootstrap} برای این کلاس رنگ تعریف کرده باشد، با ویژگی\پانویس{Specificity} برابر، آنکه دیرتر آمده برنده است؛ بنابراین ترتیب \متن‌لاتین{import} تعیین‌کننده است.
برای اطمینان بیشتر می‌توان ویژگی را بالا برد، مثلا \متن‌لاتین{ul.list-group > li.list-group-item} یا افزودن یک کلاس اختصاصی در کنار آن.

\شروع{فقرات}
\فقره ارائه‌ی کد \متن‌لاتین{CSS}: ۵۰ درصد نمره
\فقره توضیح ترتیب \متن‌لاتین{import} یا بالا بردن ویژگی: ۵۰ درصد نمره
\فقره استفاده از \متن‌لاتین{!important}: بدون نمره
\پایان{فقرات}

\end{پاسخ}
